\chapter{Barisan}\label{ch:b1:seq}

Barisan sudah dimanipulasi pada jilid Sekolah Menengah dengan gagasan
limitnya yang diterima separuh atas kepercayaan. Di sini teorinya dibangun kembali di atas
kelengkapan $\R$ (\cref{ch:b1:reals}): setiap teorema klasiknya
--- kekonvergenan monoton, \omterm{thm:b1:seq:adjacent}{barisan berdampingan}, Bolzano--Weierstrass,
kriteria Cauchy --- merupakan wajah dari aksioma tunggal itu. Bab ini
ditutup dengan telaah praktis atas barisan yang didefinisikan oleh $u_{n+1} =
f(u_n)$.

\section{Kekonvergenan}

\begin{definition}[Limit sebuah barisan]\label{def:b1:seq:limit}
Sebuah barisan $(u_n)$ berisi bilangan real \emph{konvergen} ke $\ell \in \R$
bila\index{limit!barisan}
\[
\forall \varepsilon > 0,\ \exists N \in \N,\ \forall n \geq N,
\qquad \abs{u_n - \ell} \leq \varepsilon .
\]
Kita menulis $u_n \to \ell$ atau $\lim u_n = \ell$. Barisan yang tidak
konvergen (ke bilangan real mana pun) disebut \emph{divergen}. Adapun kedivergenan \emph{ke
$+\infty$}: $\forall M,\ \exists N,\ \forall n \geq N,\ u_n \geq M$
(serupa itu untuk $-\infty$).
\end{definition}

\begin{example}[Sebuah bukti $\varepsilon$--$N$, yang ditulis lengkap sekali]\label{ex:b1:seq:epsilonN}
Klaimnya: $u_n = \dfrac{n^2 + 1}{2n^2 - 3} \to \dfrac12$. Pertama-tama
kucilkan galatnya:
\[
\Bigl| u_n - \frac12 \Bigr|
= \Bigl| \frac{2(n^2 + 1) - (2n^2 - 3)}{2(2n^2 - 3)} \Bigr|
= \frac{5}{2\,\abs{2n^2 - 3}}
= \frac{5}{2\,(2n^2 - 3)} \quad (n \geq 2).
\]
Lalu dominasikan ia dengan sesuatu yang sederhana: untuk $n \geq 2$, $2n^2 - 3
\geq n^2$, sehingga galatnya $\leq \frac{5}{2n^2} \leq \frac 5{2n}$.
Diberikan $\varepsilon > 0$, \omterm{thm:b1:reals:archimedes}{sifat Archimedes} memasok $N \geq
\max\bigl(2, \frac{5}{2\varepsilon}\bigr)$; lalu untuk $n \geq N$
galatnya $\leq \varepsilon$. Selesai. Inti gagasan penutupnya: bahwa sebuah bukti
$\varepsilon$--$N$ mempunyai tepat tiga gerakan --- menghitung
galatnya, membatasinya dengan ungkapan elementer yang menurun, lalu memecahkan
ambangnya --- dan setelah teorema pada bab ini (operasi,
apitan) orang hampir tak pernah menulis bukti semacam itu lagi: karena teoremanya
mengemas ketiga gerakan itu sekali untuk selamanya.
\end{example}

\begin{example}[Kedivergenan ke tak hingga, yang bersertifikat]\label{ex:b1:seq:toinfinity}
Klaimnya: $u_n = n^2 - 100n \to +\infty$. Faktorkan suku dominannya:
$u_n = n^2\bigl(1 - \frac{100}{n}\bigr) \geq \frac{n^2}{2}$ untuk
$n \geq 200$. Diberikan $M$, ambillah $N = \max\bigl(200,
\lceil\sqrt{2M}\rceil\bigr)$: maka untuk $n \geq N$, $u_n \geq
\frac{n^2}{2} \geq M$. Dua kebiasaan terpampang di sini: pemfaktoran suku
dominan mengubah sebuah persaingan ($n^2$ melawan $-100n$) menjadi
satu skala dikalikan faktor yang menuju $1$; sedangkan ambangnya
boleh saja raksasa (karena $u_{100} = 0$, dan barisannya bahkan negatif sebelum
$n = 100$) --- jadi kedivergenan ke $+\infty$ adalah \omterm{def:b1:logic:statement}{pernyataan} tentang
ekornya, yang acuh tak acuh terhadap sebanyak apa pun kelakuan buruk yang berhingga.
\end{example}

\begin{proposition}[Sifat yang pertama]\label{prop:b1:seq:first}
\begin{enumerate}
  \item Limitnya, bila ada, bersifat tunggal.
  \item Barisan yang konvergen bersifat terbatas.
  \item Jika $u_n \to \ell$, maka setiap perubahan atas suku yang berhingga banyak
        membiarkan kekonvergenan dan limitnya tak berubah.
\end{enumerate}
\end{proposition}

\begin{proof}
(1) Jika $u_n \to \ell$ dan $u_n \to \ell'$ dengan $\ell \neq \ell'$,
ambillah $\varepsilon = \frac{\abs{\ell - \ell'}}{3}$: maka di luar kedua
ambangnya, $\abs{\ell - \ell'} \leq \abs{\ell - u_n} + \abs{u_n -
\ell'} \leq 2\varepsilon = \frac23 \abs{\ell - \ell'}$, yang mustahil.

(2) Dengan $\varepsilon = 1$: di luar $N$, $\abs{u_n} \leq \abs\ell +
1$; sedangkan suku sebelumnya yang berhingga banyak juga terbatas, sehingga $\abs{u_n}
\leq \max(\abs{u_0}, \dots, \abs{u_{N-1}}, \abs\ell + 1)$.

(3) Secara terperinci: andaikan $v_n = u_n$ untuk $n \geq n_0$ dan $u_n \to
\ell$. Diberikan $\varepsilon > 0$, ambillah ambang $N$ bagi
$(u_n)$: maka untuk $n \geq \max(N, n_0)$, $\abs{v_n - \ell} = \abs{u_n
- \ell} \leq \varepsilon$. Jadi $v_n \to \ell$: karena definisinya
mengkuantifikasi $n \geq N$ saja, dan sebarang awalan yang hingga boleh
ditimpa dengan ongkos memperbesar ambangnya. (Inilah sebabnya
hipotesis ``untuk setiap $n$ yang besar'' mencukupi di mana-mana dalam bab
ini.)
\end{proof}

\begin{theorem}[Operasi atas limit]\label{thm:b1:seq:operations}
Jika $u_n \to \ell$ dan $v_n \to m$, maka
\[
u_n + v_n \to \ell + m, \qquad u_n v_n \to \ell m, \qquad
\frac{u_n}{v_n} \to \frac{\ell}{m} \ (\text{jika } m \neq 0),
\qquad \abs{u_n} \to \abs\ell .
\]
\end{theorem}

\begin{proof}
\emph{Jumlah:} $\abs{(u_n + v_n) - (\ell + m)} \leq \abs{u_n - \ell} +
\abs{v_n - m} \leq 2\varepsilon$ di luar ambang yang lebih besar.
\emph{Hasil kali:} tulislah
\[
u_n v_n - \ell m = (u_n - \ell)\,v_n + \ell\,(v_n - m);
\]
$(v_n)$ terbatas oleh suatu $B$ (\cref{prop:b1:seq:first}), sehingga
ruas kanannya $\leq B\abs{u_n - \ell} + \abs{\ell}\,\abs{v_n - m}$,
yang sekecil-kecilnya. \emph{Hasil bagi:} cukuplah menangani $\frac
1{v_n}$. Dengan $\varepsilon = \frac{\abs m}{2}$: di luar suatu $N_0$,
$\abs{v_n} \geq \frac{\abs m}{2}$, sehingga
\[
\Bigl| \frac{1}{v_n} - \frac 1m \Bigr|
= \frac{\abs{m - v_n}}{\abs{v_n m}}
\leq \frac{2}{m^2}\,\abs{v_n - m} \longrightarrow 0 .
\]
\emph{Nilai mutlak:} $\bigl|\abs{u_n} - \abs\ell\bigr| \leq
\abs{u_n - \ell}$ (yaitu ketaksamaan segitiga terbalik,
\cref{prop:b1:complex:rules}).
\end{proof}

\begin{example}[Operasi ditambah satu muslihat aljabar]\label{ex:b1:seq:conjugatetrick}
Hitunglah $\lim\,\bigl(\sqrt{n^2 + n} - n\bigr)$. Kedua bagiannya
secara terpisah menuju $+\infty$: sehingga teorema operasinya tak mengatakan apa pun
tentang selisihnya (yaitu sebuah \emph{bentuk tak tentu}). Kalikanlah
dengan sekawannya:
\[
\sqrt{n^2 + n} - n
= \frac{(n^2 + n) - n^2}{\sqrt{n^2 + n} + n}
= \frac{n}{\sqrt{n^2+n} + n}
= \frac{1}{\sqrt{1 + \frac1n} + 1} .
\]
Sekarang semuanya konvergen: $\sqrt{1 + \frac1n} \to 1$, karena
$0 \leq \sqrt{1 + h} - 1 = \frac{h}{\sqrt{1+h} + 1} \leq h$
(dengan sekawan lagi, lalu apitan dengan $h = \frac1n$); lalu
teorema operasinya memberikan limit $\frac{1}{1 + 1} =
\frac12$. Inti gagasan penutupnya: bahwa teorema operasi bukanlah
kalkulator bagi semua limit --- karena bentuk tak tentu ($\infty -
\infty$, $\frac00$, $0 \times \infty$, $1^\infty$) harus lebih dulu
\emph{diubah} lewat aljabar (sekawan, pemfaktoran suku
dominan) sampai setiap bagiannya konvergen; adapun mesin sistematis bagi
kasus yang membandel adalah ekspansi asimtotik pada
\cref{ch:b1:taylor}.
\end{example}

\begin{theorem}[Limit dan urutan]\label{thm:b1:seq:order}
\begin{enumerate}
  \item Jika $u_n \leq v_n$ untuk setiap $n$ yang besar, dan keduanya konvergen, maka
        $\lim u_n \leq \lim v_n$. (Adapun ketaksamaan tegas \emph{tak}
        diteruskan ke limitnya: karena $\frac 1n > 0$ tetapi $\lim = 0$.)
  \item (Teorema apit\index{teorema apit}) Jika $u_n \leq w_n
        \leq v_n$ untuk setiap $n$ yang besar dan $u_n, v_n \to \ell$, maka
        $w_n \to \ell$.
  \item Jika $u_n \to \ell > 0$, maka $u_n > \frac\ell2 > 0$ untuk setiap
        $n$ yang besar.
\end{enumerate}
\end{theorem}

\begin{proof}
(1) Andaikan $\ell = \lim u_n > m = \lim v_n$; maka dengan $\varepsilon =
\frac{\ell - m}{3}$, suku yang besar memenuhi $v_n \leq m + \varepsilon <
\ell - \varepsilon \leq u_n$, yang bertentangan dengan $u_n \leq v_n$.

(2) Di luar ambangnya: $\ell - \varepsilon \leq u_n \leq w_n \leq
v_n \leq \ell + \varepsilon$.

(3) merupakan \cref{def:b1:seq:limit} dengan $\varepsilon = \frac\ell2$.
\end{proof}

\begin{example}[Dua apitan]\label{ex:b1:seq:squeezes}
(i) $\dfrac{\sin n}{n} \to 0$: karena dari $-\frac1n \leq \frac{\sin
n}{n} \leq \frac1n$, kedua dindingnya runtuh ke $0$ --- jadi tak perlu
memahami pembilangnya yang tak beraturan itu sama sekali. (ii) Adapun $(2^n +
3^n)^{1/n} \to 3$: apitlah bagian dalamnya,
\[
3^n \leq 2^n + 3^n \leq 2\cdot3^n
\quad\Longrightarrow\quad
3 \leq (2^n + 3^n)^{1/n} \leq 3\cdot 2^{1/n} ,
\]
dan $2^{1/n} = \eu^{\frac{\ln 2}{n}} \to 1$ (seperti untuk $5^{1/n}$ pada
\cref{exo:b1:seq:2}): sehingga apitannya menghasilkan $3$. Inti gagasan
penutupnya: bahwa jumlah eksponensial yang bersaing berperilaku seperti suku yang
\emph{terbesar} --- karena yang lebih kecil terserap oleh sebuah faktor
konstanta yang tak berbahaya, yang lalu dihapus oleh akar pangkat $n$-nya.
\end{example}

\section{Barisan monoton}

\begin{theorem}[Teorema limit monoton]\label{thm:b1:seq:monotone}
Barisan naik yang terbatas di atas bersifat konvergen, ke $\sup\{u_n : n \in
\N\}$; sedangkan barisan naik yang tak terbatas di atas divergen ke
$+\infty$. (Ada \omterm{def:b1:logic:statement}{pernyataan} cerminnya untuk barisan
turun.)\index{teorema limit monoton}
\end{theorem}

\begin{proof}
Misalkan $s = \sup\{u_n\}$ (\cref{thm:b1:reals:sup}). Diberikan $\varepsilon >
0$, pencirian lewat $\varepsilon$
(\cref{prop:b1:reals:epsilon}) menghasilkan $N$ dengan $u_N > s -
\varepsilon$; lalu menurut kemonotonannya, $s - \varepsilon < u_N \leq u_n \leq
s$ untuk setiap $n \geq N$: jadi konvergen ke $s$. Jika tak terbatas: untuk setiap
$M$ ada $u_N > M$, dan kemonotonannya menjaga semua suku berikutnya di atas $M$.
\end{proof}

\begin{example}[Teorema monoton sebagai mesin keberadaan]\label{ex:b1:seq:existencemachine}
Misalkan $u_n = \prod_{k=1}^{n} \bigl(1 + \frac{1}{2^k}\bigr)$. Setiap
faktornya melampaui $1$, sehingga $(u_n)$ naik. Terbatas di atas? Ambillah
logaritmanya lalu pakai $\ln(1 + x) \leq x$
(yang didahului \cref{ex:b1:derivative:convexineq}; atau
$1 + x \leq \eu^x$ yang kasar dari jilid Sekolah Menengah):
\[
\ln u_n = \sum_{k=1}^{n} \ln\Bigl(1 + \frac{1}{2^k}\Bigr)
\leq \sum_{k=1}^{n} \frac{1}{2^k} < 1 ,
\]
sehingga $u_n < \eu$. Naik dan terbatas: jadi $(u_n)$ konvergen ke suatu
$\ell \in \intoc{u_1}{\eu}$ --- yaitu bilangan real yang terdefinisi
dengan sempurna tanpa bentuk tertutup yang terlihat ($\ell = 2.384\dots$). Inti
gagasan penutupnya: bahwa teorema limit monoton merupakan mesin keberadaan
yang paling murah dalam analisis; ia menamai $\eu$ sendiri
(\cref{ex:b1:seq:e} di bawah), dan pada \cref{ch:b1:series} ia akan
memutuskan kekonvergenan setiap deret positif lewat keterbatasan belaka.
\end{example}

\begin{theorem}[Barisan berdampingan]\label{thm:b1:seq:adjacent}
Misalkan $(a_n)$ naik dan $(b_n)$ turun, dengan $b_n - a_n \to
0$. Maka keduanya konvergen, ke limit $\ell$ yang \emph{sama}, dan $a_n
\leq \ell \leq b_n$ untuk setiap $n$.\index{barisan berdampingan}
\end{theorem}

\begin{proof}
Pertama, $a_n \leq b_n$ untuk setiap $n$: karena barisan $(b_n - a_n)$
turun dan menuju $0$, sehingga ia $\geq 0$ (sebab suku yang negatif akan
membekukannya di bawah $0$). Lalu $(a_n)$ naik dan terbatas di atas oleh
$b_0$: sehingga ia konvergen ke suatu $\ell$ (\cref{thm:b1:seq:monotone});
demikian pula $(b_n) \to \ell'$; dan $\ell' - \ell = \lim (b_n - a_n) =
0$. Adapun ketaksamaan $a_n \leq \ell \leq b_n$ menyusul dari
kemonotonannya (yaitu $\ell = \sup a_k \geq a_n$, dan seterusnya).
\end{proof}

\begin{omfigure}
\begin{tikzpicture}[scale=10.5]
  \draw[-Stealth] (-0.03,0) -- (1.08,0);
  % a_n increasing (omDef), b_n decreasing (omProp)
  \fill[omDef] (0,0) circle (0.3pt) node[below=2pt, font=\small]
    {$a_0$};
  \fill[omDef] (0.18,0) circle (0.3pt) node[below=2pt,
    font=\small] {$a_1$};
  \fill[omDef] (0.30,0) circle (0.3pt) node[below=2pt,
    font=\small] {$a_2$};
  \fill[omDef] (0.37,0) circle (0.3pt);
  \fill[omDef] (0.41,0) circle (0.3pt);
  \fill[omProp] (1.0,0) circle (0.3pt) node[below=2pt,
    font=\small] {$b_0$};
  \fill[omProp] (0.78,0) circle (0.3pt) node[below=2pt,
    font=\small] {$b_1$};
  \fill[omProp] (0.62,0) circle (0.3pt) node[below=2pt,
    font=\small] {$b_2$};
  \fill[omProp] (0.54,0) circle (0.3pt);
  \fill[omProp] (0.50,0) circle (0.3pt);
  \draw[omThm, thick] (0.455,-0.014) -- (0.455,0.014);
  \node[above, omThm, font=\small] at (0.455,0.02) {$\ell$};
  \draw[Stealth-Stealth, gray] (0.30,0.05) -- (0.62,0.05);
  \node[above, font=\small, text=gray] at (0.46,0.055)
    {$b_n - a_n \to 0$};
\end{tikzpicture}

{\small \omterm{thm:b1:seq:adjacent}{Barisan berdampingan}: di sini $(a_n)$ mendaki, $(b_n)$ menurun,
dan celah di antaranya menyusut ke $0$. Setiap \omterm{prop:b1:reals:intervals}{selang}
$\intcc{a_n}{b_n}$ memuat semua \omterm{prop:b1:reals:intervals}{selang} berikutnya, dan limit bersamanya
$\ell$ adalah satu-satunya titik yang tersisa pada setiap \omterm{prop:b1:reals:intervals}{selang} --- yaitu
gambaran di balik bukti dikotomi Bolzano--Weierstrass
di bawah dan bukti teorema nilai antara pada
\cref{ch:b1:continuity}.}
\end{omfigure}

\begin{example}[Bilangan $\eu$]\label{ex:b1:seq:e}
Tetapkan $a_n = \sum_{k=0}^{n} \frac{1}{k!}$ dan $b_n = a_n + \frac{1}{n
\cdot n!}$ (dengan $n \geq 1$). Maka $(a_n)$ naik; dan
\[
b_{n+1} - b_n = \frac{1}{(n+1)!} + \frac{1}{(n+1)(n+1)!} -
\frac{1}{n\,n!}
= \frac{n(n+1) + n - (n+1)^2}{n(n+1)(n+1)!}
= \frac{-1}{n(n+1)(n+1)!} < 0 ,
\]
sehingga $(b_n)$ turun, dan $b_n - a_n \to 0$: jadi berdampingan. Limit
bersamanya adalah (menurut definisi di sini) bilangan $\eu \approx 2.71828$; dan
ketaksamaan $a_n < \eu < b_n$ cukup tajam untuk membuktikan $\eu \notin
\Q$ (\cref{exo:b1:seq:9}).
\end{example}

\section{Subbarisan dan Bolzano--Weierstrass}

\begin{definition}[Subbarisan]\label{def:b1:seq:subsequence}
Sebuah \emph{subbarisan}\index{subbarisan} dari $(u_n)$ adalah barisan
$(u_{\varphi(n)})$ dengan $\varphi \colon \N \to \N$ yang naik
tegas (perhatikan $\varphi(n) \geq n$, lewat induksi).
\end{definition}

\begin{proposition}\label{prop:b1:seq:subsequences}
Jika $u_n \to \ell$ (dengan $\ell \in \R$ atau $\pm\infty$), maka setiap \omterm{def:b1:seq:subsequence}{subbarisannya}
menuju $\ell$. Akibatnya, barisan yang mempunyai dua \omterm{def:b1:seq:subsequence}{subbarisan} dengan
limit yang berbeda bersifat divergen. Sebaliknya, jika $(u_{2n})$ dan
$(u_{2n+1})$ sama-sama konvergen ke $\ell$ yang \emph{sama}, maka $u_n \to
\ell$.
\end{proposition}

\begin{proof}
Di luar ambang $N$ bagi $(u_n)$, semua indeks $\varphi(n) \geq n
\geq N$ memenuhi syaratnya (adapun ketaksamaan $\varphi(n) \geq n$ merupakan
induksi yang dicatat pada \cref{def:b1:seq:subsequence}: $\varphi(0)
\geq 0$, dan $\varphi(n+1) > \varphi(n) \geq n$ memaksa
$\varphi(n+1) \geq n + 1$). Untuk konversnya: diberikan
$\varepsilon$, ambillah kedua ambangnya $N_0$ (yang genap) dan $N_1$
(yang ganjil); maka sebarang indeks $n \geq \max(2N_0, 2N_1 + 1)$
entah genap, $n = 2k$ dengan $k \geq N_0$, entah ganjil, $n = 2k+1$ dengan
$k \geq N_1$ --- dan pada kedua kasusnya $\abs{u_n - \ell} \leq
\varepsilon$: jadi setiap indeksnya tertutupi oleh salah satu dari kedua
\omterm{def:b1:seq:subsequence}{subbarisan} itu, dan itulah seluruh intinya.
\end{proof}

\begin{example}[Limit subbarisan]\label{ex:b1:seq:sublimits}
Untuk $u_n = (-1)^n \frac{n}{n+1}$: \omterm{def:b1:seq:subsequence}{subbarisan} genapnya menuju
$1$, sedangkan yang ganjil ke $-1$, sehingga barisannya divergen --- tetapi ia
berbuat demikian secara tertata, dengan menggerombol di sekitar kedua nilai $\pm 1$.
Untuk $u_n = \cos\frac{2\pi n}{3}$: ketiga \omterm{def:b1:seq:subsequence}{subbarisan}
berindeks $3k$, $3k + 1$, $3k + 2$ bersifat konstan, yang bernilai $1$,
$-\frac12$, $-\frac12$; jadi \omterm{def:b1:logic:sets}{himpunan} limit \omterm{def:b1:seq:subsequence}{subbarisannya} adalah $\{1,
-\frac12\}$. Inti gagasan penutupnya: bahwa barisan yang \emph{terbatas}
konvergen tepat ketika ia hanya mempunyai satu limit \omterm{def:b1:seq:subsequence}{subbarisan}
(\cref{exo:b1:seq:8}); jadi kedivergenan barisan yang terbatas selalu
berarti sekurang-kurangnya dua gerombolan, dan Bolzano--Weierstrass di bawah
menjamin bahwa ada sekurang-kurangnya satu.
\end{example}

\begin{theorem}[Bolzano--Weierstrass]\label{thm:b1:seq:bw}
Setiap barisan bilangan real yang terbatas mempunyai \omterm{def:b1:seq:subsequence}{subbarisan} yang
konvergen.\index{teorema Bolzano--Weierstrass}
\end{theorem}

\begin{proof}
Misalkan $u_n \in \intcc{a}{b}$ untuk setiap $n$. Bangunlah ruas yang bersarang lewat
dikotomi: tetapkan $\intcc{a_0}{b_0} = \intcc{a}{b}$; lalu diberikan
$\intcc{a_k}{b_k}$ yang memuat $u_n$ untuk tak hingga banyak $n$, salah satu
dari kedua paruhnya masih memuat $u_n$ untuk tak hingga banyak $n$ --- namailah
ia $\intcc{a_{k+1}}{b_{k+1}}$. Barisan $(a_k)$ dan $(b_k)$ itu
berdampingan (karena $b_k - a_k = \frac{b-a}{2^k} \to 0$), dengan limit bersama
$\ell$ (\cref{thm:b1:seq:adjacent}).

Ekstraksinya: pilihlah $\varphi(0)$ dengan $u_{\varphi(0)} \in
\intcc{a_0}{b_0}$, lalu, secara induktif, $\varphi(k+1) > \varphi(k)$
dengan $u_{\varphi(k+1)} \in \intcc{a_{k+1}}{b_{k+1}}$ --- yang mungkin
karena ruas itu memuat tak hingga banyak sukunya. Maka $a_k \leq
u_{\varphi(k)} \leq b_k$, dan teorema apitnya memberikan
$u_{\varphi(k)} \to \ell$.
\end{proof}

\begin{remark}[Yang dikatakan Bolzano--Weierstrass, dan yang tidak]
Ia \emph{memang} mengatakan: bahwa dari keterbatasan belaka, suatu \omterm{def:b1:seq:subsequence}{subbarisan}
konvergen --- yaitu keberadaan tanpa rumus, sebagaimana dijelaskan bukti
dikotominya (karena tak ada yang memberitahu kita indeks \emph{mana} yang bertahan). Ia
\emph{tak} mengatakan bahwa limitnya tunggal: karena $((-1)^n)$ mempunyai
\omterm{def:b1:seq:subsequence}{subbarisan} yang konvergen ke $1$ dan ke $-1$, dan \omterm{def:b1:logic:sets}{himpunan}
limit \omterm{def:b1:seq:subsequence}{subbarisannya} bahkan boleh tak hingga
(\cref{ex:b1:seq:sublimits}, beserta seluruh \omterm{def:b1:logic:sets}{himpunan} Cantor pada
\cref{pb:b1:topology:1}). Ia tak bertahan pada ketakterbatasan:
karena $(n)$ sama sekali tak mempunyai \omterm{def:b1:seq:subsequence}{subbarisan} yang konvergen --- meskipun kita
selalu dapat mengekstrak \omterm{def:b1:seq:subsequence}{subbarisan} yang menuju $+\infty$ atau $-\infty$
dari sebarang barisan yang tak terbatas (misalnya dengan memilih $\varphi(k)$ dengan
$u_{\varphi(k)} \geq k$). Bila dipakai dengan benar, teoremanya menjadi sebuah
\emph{pompa keberadaan}: ia muncul pada titik genting kriteria
Cauchy di bawah, pada teorema Heine, dan pada teorema nilai
ekstrem --- selalu untuk menghasilkan sebuah titik yang tak disodorkan konstruksi
eksplisit mana pun.
\end{remark}

\section{Barisan Cauchy dan kelengkapan}

\begin{definition}\label{def:b1:seq:cauchy}
Sebuah barisan $(u_n)$ disebut \emph{barisan Cauchy}\index{barisan Cauchy}
bila sukunya menjadi sedekat-dekatnya \emph{satu sama lain}:
\[
\forall \varepsilon > 0,\ \exists N,\ \forall p, q \geq N,
\qquad \abs{u_p - u_q} \leq \varepsilon .
\]
\end{definition}

\begin{example}[Memeriksa sifat Cauchy dengan tangan]\label{ex:b1:seq:cauchyhand}
Misalkan $u_n = \sum_{k=0}^{n} \frac{\cos k}{2^k}$ --- tanpa kemonotonan,
tanpa limit yang dapat ditebak. Untuk $p > q$:
\[
\abs{u_p - u_q}
= \Bigl| \sum_{k=q+1}^{p} \frac{\cos k}{2^k} \Bigr|
\leq \sum_{k=q+1}^{p} \frac{1}{2^k}
< \frac{1}{2^{q}} ,
\]
menurut ketaksamaan segitiga, $\abs{\cos k} \leq 1$ dan jumlah geometri
yang hingga. Diberikan $\varepsilon > 0$, pilihlah $N$ dengan $2^{-N}
\leq \varepsilon$: maka semua celah di luar $N$ bernilai $\leq \varepsilon$,
sehingga barisannya Cauchy, jadi konvergen --- ke limit yang tak seorang pun
dapat menamainya dalam bentuk tertutup, dan itu justru intinya. Inti gagasan
penutupnya: bahwa dominasi geometri atas pertambahannya yang menjadi cara baku
untuk memperoleh sifat Cauchy, dan \cref{ch:b1:series} akan
membotolkan argumen itu sebagai ``kekonvergenan mutlak mengakibatkan
kekonvergenan''.
\end{example}

\begin{theorem}[Kelengkapan $\R$]\label{thm:b1:seq:complete}
Barisan bilangan real konvergen jika dan hanya jika ia \omterm{def:b1:seq:cauchy}{barisan
Cauchy}.\index{kelengkapan R@kelengkapan $\R$}
\end{theorem}

\begin{proof}
($\Rightarrow$) Jika $u_n \to \ell$: maka di luar ambang bagi
$\frac\varepsilon2$, $\abs{u_p - u_q} \leq \abs{u_p - \ell} +
\abs{\ell - u_q} \leq \varepsilon$.

($\Leftarrow$) Misalkan $(u_n)$ Cauchy. \emph{Ia terbatas}: karena dengan
$\varepsilon = 1$, di luar $N$ semua sukunya terletak sejauh $1$ dari $u_N$, dan
kepalanya berhingga. \emph{Ekstraksinya}: menurut
\cref{thm:b1:seq:bw}, suatu \omterm{def:b1:seq:subsequence}{subbarisan} $u_{\varphi(n)} \to \ell$.
\emph{Kesimpulannya}: diberikan $\varepsilon > 0$, ambillah $N$ (dari Cauchy, untuk
$\frac\varepsilon2$) dan $n \geq N$ dengan $\abs{u_{\varphi(n)} - \ell}
\leq \frac\varepsilon2$ dan $\varphi(n) \geq N$; maka untuk setiap $p
\geq N$:
\[
\abs{u_p - \ell} \leq \abs{u_p - u_{\varphi(n)}} +
\abs{u_{\varphi(n)} - \ell} \leq \varepsilon . \qedhere
\]
\end{proof}

\begin{remark}
Nilai kriterianya: bahwa ia mengesahkan kekonvergenan \emph{tanpa
menamai limitnya}. Ia gagal di atas $\Q$ (karena pemenggalan desimal
$\sqrt 2$ membentuk \omterm{def:b1:seq:cauchy}{barisan Cauchy} berisi bilangan rasional tanpa limit yang
rasional): jadi kelengkapan merupakan sifat $\R$, yang setara dengan aksioma
\omterm{def:b1:reals:bounds}{batas atas}. Ia juga kuda beban di balik kekonvergenan
deret (\cref{ch:b1:series}).
\end{remark}

\begin{example}[Barisan Cauchy yang limitnya tak kasatmata]\label{ex:b1:seq:basel}
Misalkan $S_n = \sum_{k=1}^{n} \frac{1}{k^2}$. Untuk $p > q \geq 1$:
\[
S_p - S_q = \sum_{k=q+1}^{p} \frac{1}{k^2}
\leq \sum_{k=q+1}^{p} \frac{1}{k(k-1)}
= \sum_{k=q+1}^{p} \Bigl(\frac{1}{k-1} - \frac 1k\Bigr)
= \frac 1q - \frac 1p < \frac 1q ,
\]
sehingga di luar $N > \frac1\varepsilon$ semua celahnya $\leq \varepsilon$:
jadi $(S_n)$ Cauchy, sehingga konvergen. Perhatikan apa yang baru saja terjadi: kita
membuktikan bahwa sebuah bilangan real tertentu ada tanpa mempunyai nama apa pun
untuknya. (Ia $\frac{\pi^2}{6}$ --- yaitu kesamaan Euler yang termasyhur,
yang dibuktikan pada jilid Tahun ke-2; dan tak ada apa pun dalam bab ini yang dapat
memberitahu kita hal itu.) Pembagian kerja inilah --- keberadaan sekarang,
pengenalan belakangan, kalau pun pernah --- yang menjadi seluruh inti kriteria
Cauchy, dan mesin teori deret pada
\cref{ch:b1:series}.
\end{example}

\section{Barisan rekuren}

\begin{method}[Menelaah $u_{n+1} = f(u_n)$]\label{met:b1:seq:recurrent}
Diberikan $f$ dan titik awal $u_0$:\index{barisan rekuren}
\begin{enumerate}
  \item \emph{\omterm{prop:b1:reals:intervals}{Selang} yang stabil:} carilah \omterm{prop:b1:reals:intervals}{selang} $I$ dengan $f(I)
        \subseteq I$ yang memuat $u_0$: maka semua $u_n \in I$ (lewat
        induksi).
  \item \emph{Calon limitnya:} jika $u_n \to \ell \in I$ dan $f$
        kontinu di $\ell$ (\cref{ch:b1:continuity}), maka $\ell$
        merupakan \emph{titik tetap}: $f(\ell) = \ell$. Pecahkanlah $f(x) = x$.
  \item \emph{Kemonotonannya:} jika $f$ naik pada $I$, maka
        $(u_n)$ monoton (naik bila $u_1 \geq u_0$, dan
        turun bila tidak); lalu digabung dengan keterbatasan,
        \cref{thm:b1:seq:monotone} menyimpulkannya. Jika $f$ turun,
        telaahlah kedua \omterm{def:b1:seq:subsequence}{subbarisan} $(u_{2n})$ dan $(u_{2n+1})$,
        yang monoton bagi $f \circ f$.
  \item \emph{Kendali galatnya:} sebuah ketaksamaan $\abs{f(x) - \ell} \leq
        k\abs{x - \ell}$ dengan $k < 1$ memberikan $\abs{u_n - \ell} \leq
        k^n \abs{u_0 - \ell} \to 0$ secara langsung.
\end{enumerate}
\end{method}

\begin{example}[Metode Heron]\label{ex:b1:seq:heron}
Misalkan $u_0 = 2$ dan $u_{n+1} = \dfrac12\Bigl(u_n +
\dfrac{2}{u_n}\Bigr)$: yaitu algoritma purba bagi
$\sqrt 2$.\index{metode Heron}
\begin{itemize}
  \item \emph{Kestabilannya:} untuk $x > 0$, ketaksamaan rata-rata
        aritmetika--geometri memberikan $\frac12(x + \frac2x) \geq \sqrt{x \cdot
        \frac 2x} = \sqrt 2$; sehingga $I = \intco{\sqrt 2}{+\infty}$
        stabil dan memuat $u_1$ (memang $u_1 = \frac32 \geq
        \sqrt 2$).
  \item \emph{Kemonotonannya:} untuk $x \geq \sqrt 2$, $\;x - f(x) =
        \frac{x^2 - 2}{2x} \geq 0$: sehingga barisannya turun mulai
        $u_1$, dan terbatas di bawah oleh $\sqrt 2$: jadi ia konvergen.
  \item \emph{Limitnya:} titik tetapnya memecahkan $x = \frac12(x +
        \frac2x)$, yakni $x^2 = 2$: jadi pada $I$, $\ell = \sqrt 2$.
  \item \emph{Kecepatannya:} $u_{n+1} - \sqrt 2 = \frac{(u_n -
        \sqrt2)^2}{2u_n} \leq \frac{(u_n - \sqrt 2)^2}{2\sqrt 2}$:
        sehingga banyaknya angka yang benar kira-kira \emph{berlipat dua} pada setiap
        langkahnya (yaitu kekonvergenan kuadratik).
\end{itemize}
\end{example}

\begin{omfigure}
\begin{tikzpicture}
\begin{axis}[omaxis, width=8.6cm, height=7cm,
    xmin=0.9, xmax=2.35, ymin=0.9, ymax=2.35,
    xtick={1,1.5,2}, ytick={1,1.5,2},
    xlabel={$x$}, ylabel={$y$}]
  \addplot[gray, dashed, domain=0.9:2.3] {x}
    node[pos=0.72, above left, font=\small] {$y = x$};
  \addplot[omDef, very thick, domain=0.95:2.3, samples=100]
    {0.5*(x + 2/x)} node[pos=0.88, above left, font=\small]
    {$y = f(x)$};
  % cobweb: u0=2, u1=1.5, u2=1.4167
  \draw[gray, dashed, thin] (axis cs:1.5,0.9) -- (axis cs:1.5,1.4167);
  \draw[omThm, thick, -Stealth] (axis cs:2,0.9) -- (axis cs:2,1.5);
  \draw[omThm, thick, -Stealth] (axis cs:2,1.5) -- (axis cs:1.5,1.5);
  \draw[omThm, thick, -Stealth] (axis cs:1.5,1.5) --
    (axis cs:1.5,1.4167);
  \draw[omThm, thick, -Stealth] (axis cs:1.5,1.4167) --
    (axis cs:1.4167,1.4167);
  \fill (axis cs:1.41421,1.41421) circle (1.6pt);
  \node[anchor=north, font=\small] at (axis cs:1.4142,1.385)
    {$\sqrt 2$};
  \node[anchor=west, font=\small, omThm] at (axis cs:2.02,1.05)
    {$u_0$};
  \node[anchor=east, font=\small, omThm] at (axis cs:1.48,1.05)
    {$u_1$};
\end{axis}
\end{tikzpicture}

{\small Iterasi Heron $u_{n+1} = \frac12\bigl(u_n +
\frac{2}{u_n}\bigr)$, yang digambar sebagai tangga di antara grafik $f$
dan diagonal $y = x$: dari $u_0 = 2$, iterasinya meluncur turun ke
titik tetap $\sqrt 2$.}
\end{omfigure}

\begin{remark}[Jebakan yang lazim dengan limit]
Empat yang klasik. (i) \emph{Langkah yang kecil tak mengakibatkan kekonvergenan}:
karena $u_{n+1} - u_n \to 0$ jauh lebih lemah daripada sifat Cauchy ---
jumlah harmonik $H_n$ berlangkah $\frac{1}{n+1} \to 0$ namun
divergen ke $+\infty$ (\cref{exo:b1:seq:5}); sebab syarat Cauchy
mengendalikan $\abs{u_p - u_q}$ bagi \emph{semua} pasangan yang besar, bukan
yang berurutan. (ii) \emph{Ketaksamaan tegas mati pada
limitnya}: dari $u_n < v_n$ untuk setiap $n$ kita hanya memperoleh $\lim u_n
\leq \lim v_n$ (\cref{thm:b1:seq:order}); karena $\frac1n > 0$ namun
$\lim = 0$. (iii) \emph{Terbatas bukan berarti konvergen}: karena $((-1)^n)$
terbatas dan divergen; jadi keterbatasan ditambah \emph{kemonotonan}
menghasilkan kekonvergenan, sedangkan keterbatasan belaka hanya menjamin \omterm{def:b1:seq:subsequence}{subbarisan}
yang konvergen (\cref{thm:b1:seq:bw}). (iv) \emph{Persamaan titik
tetap datang kedua, bukan pertama}: karena untuk $u_{n+1} = f(u_n)$,
memecahkan $f(\ell) = \ell$ mengenali limitnya \emph{hanya setelah}
kekonvergenannya terbukti. Rekurensi $u_{n+1} = 2u_n$ mempunyai
titik tetap tunggal $\ell = 0$, namun dari $u_0 = 1$ barisannya
melesat ke $+\infty$: jadi persamaan $\ell = 2\ell$ tak pernah
berhak atas sebuah limit. Urutan pengerjaannya, selalu: keberadaan
dulu (\cref{met:b1:seq:recurrent}, langkah 1--3), lalu pengenalan
kedua.
\end{remark}

\begin{example}[Sebuah $f$ yang turun: rekurensi emas]\label{ex:b1:seq:golden}
Misalkan $u_0 = 1$ dan $u_{n+1} = \dfrac{1}{1 + u_n}$. Di sini $f(x) =
\frac{1}{1+x}$ bersifat \emph{turun}, sehingga barisannya tak
monoton (ia berselang-seling di sekitar limitnya); jadi langkah kontraksi
pada \cref{met:b1:seq:recurrent} yang menjadi perkakas yang tepat. Kestabilannya: jika
$x \in \intcc{\frac12}{1}$ maka $1 + x \in
\intcc{\frac32}{2}$, sehingga $f(x) \in \intcc{\frac12}{\frac23}
\subseteq \intcc{\frac12}{1}$, dan $u_1 = \frac12$ menaruh
seluruh barisannya di sana. Titik tetapnya: $\ell = \frac{1}{1+\ell}$
dengan $\ell > 0$ memberikan $\ell^2 + \ell - 1 = 0$, yakni
\[
\ell = \frac{\sqrt5 - 1}{2} = 0.6180\dots
\]
(yaitu invers rasio emas). Kontraksinya: untuk $x, y \in
\intcc{\frac12}{1}$,
\[
\abs{f(x) - f(y)} = \frac{\abs{x - y}}{(1+x)(1+y)}
\leq \frac{\abs{x-y}}{(3/2)^2} = \frac49\,\abs{x - y} ,
\]
sehingga $\abs{u_n - \ell} \leq \bigl(\frac49\bigr)^{n-1}\abs{u_1 -
\ell} \to 0$: jadi konvergen, dengan kecepatan geometri, tanpa
memerlukan kemonotonan. Inti gagasan penutupnya: bahwa metode monoton dan
metode kontraksi \omterm{def:b1:arith:divides}{membagi} dunia rekuren di antara keduanya ---
karena $f$ yang naik memberikan orbit yang monoton, sedangkan $f$ yang turun memberikan
orbit berselang-seling yang dijinakkan oleh konstanta Lipschitz $< 1$ (adapun
teori sistematisnya adalah \cref{exo:b1:derivative:11}).
\end{example}

\begin{remark}[Cakrawala di dalam jilid ini]
Barisan adalah alat ukur yang ditodongkan sisa jilid ini
kepada setiap objek. Pada \cref{ch:b1:topology} ia
\emph{mencirikan} ketertutupan dan kekompakan; pada
\cref{ch:b1:continuity} ia mengangkut limit fungsi; pada
\cref{ch:b1:integration} jumlah Riemann merupakan barisan yang konvergen
ke integralnya; sedangkan \cref{ch:b1:series} \emph{adalah} teori satu
kelas khusus barisan, yaitu jumlah parsialnya. Bahkan bab aljabarnya
memakainya: iterasi sebuah matriks pada
\cref{ch:b1:matrices} membentuk barisan yang perilakunya
(kekonvergenan $A^n$) merupakan pertanyaan aljabar linear dengan
kosakata bab ini. Adapun dua teorema yang harus dibawa ke mana-mana:
limit monoton (yaitu keberadaan dari urutan) dan Bolzano--Weierstrass
(keberadaan dari keterbatasan) --- di antara keduanya, hampir setiap
limit dalam buku ini terlahir.
\end{remark}

\begin{remark}[Barisan kompleks]
Barisan $(z_n)$ berisi bilangan kompleks konvergen ke $\ell$ bila
$\abs{z_n - \ell} \to 0$; setara dengan itu, bila $\Re(z_n) \to \Re(\ell)$
dan $\Im(z_n) \to \Im(\ell)$ (bandingkan $\abs{z}$ dengan $\abs{\Re z} +
\abs{\Im z}$). Adapun teorema yang tak melibatkan urutan --- operasi,
Bolzano--Weierstrass (dengan mengekstrak dua kali), kriteria Cauchy --- terbawa
persis apa adanya.
\end{remark}

\section{Latihan}

\begin{exercise}[$\star$]\label{exo:b1:seq:1}
Langsung dari \cref{def:b1:seq:limit}, buktikan bahwa
$\dfrac{2n+1}{n+3} \to 2$, dan bahwa $(u_n) = ((-1)^n)$ divergen.
\end{exercise}

\begin{exercise}[$\star$]\label{exo:b1:seq:2}
Hitunglah limitnya:
\[
\frac{n^2 - 3n + 1}{2n^2 + 5},
\qquad
\sqrt{n+1} - \sqrt n,
\qquad
\frac{2^n + n^3}{3^n - n^2},
\qquad
\sqrt[n]{5}\ \Bigl(= 5^{1/n}\Bigr).
\]
\end{exercise}

\begin{exercise}[$\star$]\label{exo:b1:seq:3}
Buktikan perbandingan yang baku: jika $\abs{q} < 1$ maka $q^n \to 0$
\emph{(tulislah $\frac{1}{\abs q} = 1 + h$, $h > 0$, lalu pakai
ketaksamaan Bernoulli $(1+h)^n \geq 1 + nh$, yang dibuktikan lewat
induksi)}. Bagaimana perilakunya untuk $q = 1$, $q = -1$, $\abs q >
1$?
\end{exercise}

\begin{exercise}[$\star$]\label{exo:b1:seq:4}
Misalkan $u_{n+1} = \frac{u_n + 3}{2}$, $u_0 = 0$. Carilah titik tetapnya
$\ell$, buktikan bahwa $v_n = u_n - \ell$ bersifat geometri, lalu berikan rumus
eksplisit beserta limit $(u_n)$.
\end{exercise}

\begin{exercise}[$\star\star$]\label{exo:b1:seq:5}
(Deret harmonik) Misalkan $H_n = \sum_{k=1}^{n} \frac 1k$. Buktikan bahwa
$H_{2n} - H_n \geq \frac12$ untuk setiap $n \geq 1$, lalu simpulkan bahwa
$(H_n)$ \emph{bukan} \omterm{def:b1:seq:cauchy}{barisan Cauchy}, sehingga divergen (ke
$+\infty$, karena ia naik).
\end{exercise}

\begin{exercise}[$\star\star$]\label{exo:b1:seq:6}
Andaikan $(u_{2n})$, $(u_{2n+1})$ dan $(u_{3n})$ semuanya konvergen. Buktikan
bahwa $(u_n)$ konvergen. \emph{(Carilah \omterm{def:b1:seq:subsequence}{subbarisan} bersama untuk menyamakan
limitnya.)}
\end{exercise}

\begin{exercise}[$\star\star$]\label{exo:b1:seq:7}
Telaahlah barisan $u_0 = 0$, $u_{n+1} = \sqrt{2 + u_n}$: yaitu
kestabilan, kemonotonan, dan limitnya. Lalu buktikan batas galatnya
$\abs{u_n - 2} \leq \dfrac{2}{3^{\,n}}$ \emph{(tunjukkan $2 - u_{n+1} =
\dfrac{2 - u_n}{2 + \sqrt{2 + u_n}}$ lalu batasi penyebutnya dari bawah
oleh $3$)}.
\end{exercise}

\begin{exercise}[$\star\star$]\label{exo:b1:seq:8}
Misalkan $(u_n)$ terbatas, sedemikian sehingga setiap \omterm{def:b1:seq:subsequence}{subbarisan}
$(u_n)$ yang konvergen mempunyai limit $\ell$ yang \emph{sama}. Buktikan $u_n \to \ell$.
\emph{(Pertentangan ditambah Bolzano--Weierstrass.)}
\end{exercise}

\begin{exercise}[$\star\star\star$]\label{exo:b1:seq:9}
Dengan notasi \cref{ex:b1:seq:e}, andaikan $\eu = \frac pq$
dengan $p, q \in \N^*$. Dengan memakai $a_q < \eu < b_q = a_q + \frac{1}{q\,
q!}$, kalikanlah dengan $q!$ lalu turunkan sebuah pertentangan antara dua
bilangan bulat. Simpulkan: $\eu$ irasional.
\end{exercise}

\begin{exercise}[$\star\star\star$]\label{exo:b1:seq:10}
(Rata-rata Cesàro) Untuk barisan $(u_n)_{n \geq 1}$, tetapkan $c_n =
\frac{u_1 + \dots + u_n}{n}$.
\begin{enumerate}
  \item Buktikan bahwa $u_n \to \ell$ mengakibatkan $c_n \to \ell$ \emph{(potonglah
        jumlahnya pada sebuah ambang $N$; lalu batasi kepalanya oleh besaran
        tetap dibagi $n$, dan ekornya oleh $\varepsilon$)}.
  \item Tunjukkan lewat contoh bahwa konversnya gagal.
  \item Simpulkan bahwa jika $u_{n+1} - u_n \to \ell$, maka $\frac{u_n}{n}
        \to \ell$.
\end{enumerate}
\end{exercise}

\begin{exercise}[$\star\star\star$]\label{exo:b1:seq:11}
Misalkan $(u_n)$ memenuhi $0 \leq u_{m+n} \leq u_m + u_n$ untuk setiap $m, n$
(yaitu subaditif). Buktikan bahwa $\bigl(\frac{u_n}{n}\bigr)$ konvergen ke
$\inf_{n \geq 1} \frac{u_n}{n}$. \emph{(Untuk $m$ yang tetap, tulislah $n = qm
+ r$ lalu batasi $\frac{u_n}{n}$ memakai $u_n \leq q\,u_m + u_r$.)}
\end{exercise}

\begin{exercise}[$\star\star\star$]\label{exo:b1:seq:12}
Dengan memakai kepadatan \omterm{def:b1:structures:subgroup}{subgrup} $\Z + 2\pi\Z$ dari $(\R, +)$
(\cref{exo:b1:reals:9}), buktikan bahwa barisan $(\sin n)_{n \in
\N}$ padat di $\intcc{-1}{1}$ --- khususnya ia divergen.
\end{exercise}

\section{Soal: Cesàro, Stolz, dan jatuhnya sinus yang lambat}

\begin{problem}[{Soal akhir pekan --- teorema Cesàro--Stolz
dan asimtotik $u_n \sim \sqrt{3/n}$ untuk $u_{n+1} = \sin
u_n$}]
\label{pb:b1:seq:1}
Teorema Cesàro--Stolz adalah aturan l'Hospital yang diskret: untuk
mencari limit sebuah hasil bagi $a_n/b_n$, cukuplah mencari
limit hasil bagi \emph{selisihnya} $(a_{n+1} -
a_n)/(b_{n+1} - b_n)$. Soal ini membuktikan teoremanya, memanen
limit klasik dengannya, lalu membidikkannya pada sasaran yang termasyhur:
yaitu barisan $u_{n+1} = \sin u_n$, yang merayap ke $0$ dengan
kecepatan yang terhitung persis $u_n \sim \sqrt{3/n}$.\index{teorema
Cesaro--Stolz@teorema Cesàro--Stolz} Dua fakta dari jilid Sekolah
Menengah diberikan cuma-cuma di sini dan dibuktikan kembali dengan jujur nanti dalam jilid
ini: yaitu ketaksamaan garis singgungnya
\[
\tag{G1} \eu^{u} \geq 1 + u \quad (u \in \R),
\]
yang dibuktikan kembali lewat kecembungan pada \cref{ch:b1:derivative}, dan apitan
sinusnya
\[
\tag{G2} x - \frac{x^3}{6} \;\leq\; \sin x \;\leq\; x -
\frac{x^3}{6} + \frac{x^5}{120} \quad (0 \leq x \leq 1), \qquad
\abs{\sin x} \leq \abs{x} \quad (x \in \R),
\]
yang dibuktikan kembali lewat rumus Taylor pada \cref{ch:b1:taylor}.

\medskip
\noindent\textbf{Bagian I --- Jumlah tanpa rumus tertutup.}
\begin{enumerate}
  \item Dengan memakai $1 + 2 + \dots + n = \frac{n(n+1)}{2}$ dan $1^2 +
        \dots + n^2 = \frac{n(n+1)(2n+1)}{6}$, hitunglah
        $\lim \frac{1 + 2 + \dots + n}{n^2}$ dan $\lim \frac{1^2 +
        \dots + n^2}{n^3}$.
  \item Misalkan $T_n = \sum_{k=1}^n \sqrt k$, yang tak mempunyai rumus
        tertutup. Buktikan apitan
        \[
        \frac{1}{2\sqrt 2}\,n^{3/2} \;\leq\; T_n \;\leq\;
        n^{3/2}
        \]
        \emph{(pertahankan hanya suku dengan $k > \frac n2$ untuk batas
        bawahnya)}. Jadi $T_n$ berorde $n^{3/2}$ --- tetapi dengan
        konstanta yang mana? Tahanlah pertanyaannya sampai pertanyaan 8.
  \item (Lema teleskopis) Misalkan $(b_n)$ naik tegas
        dan andaikan bahwa untuk setiap $k \geq N$,
        \[
        m \;\leq\; \frac{a_{k+1} - a_k}{b_{k+1} - b_k} \;\leq\;
        M .
        \]
        Buktikan bahwa $m \leq \dfrac{a_n - a_N}{b_n - b_N} \leq M$
        untuk setiap $n > N$.
\end{enumerate}

\medskip
\noindent\textbf{Bagian II --- Teorema Cesàro--Stolz.} Misalkan
$(b_n)$ naik tegas dengan $b_n \to +\infty$, dan
andaikan $\dfrac{a_{n+1} - a_n}{b_{n+1} - b_n} \to \ell \in \R$.
\begin{enumerate}[resume]
  \item Tetapkan $\varepsilon > 0$. Tunjukkan bahwa ada $N$ sedemikian sehingga
        $\ell - \varepsilon \leq \dfrac{a_n - a_N}{b_n - b_N}
        \leq \ell + \varepsilon$ untuk setiap $n > N$.
  \item Tegakkan, untuk $n > N$, kesamaan
        \[
        \frac{a_n}{b_n} - \ell
        = \frac{a_N - \ell\,b_N}{b_n}
        + \Bigl(1 - \frac{b_N}{b_n}\Bigr)
        \Bigl(\frac{a_n - a_N}{b_n - b_N} - \ell\Bigr),
        \]
        lalu simpulkan teoremanya: $\dfrac{a_n}{b_n} \to \ell$.
  \item Buktikan varian $+\infty$-nya: bahwa jika $\dfrac{a_{n+1} -
        a_n}{b_{n+1} - b_n} \to +\infty$ (dengan hipotesis yang sama atas
        $(b_n)$), maka $\dfrac{a_n}{b_n} \to +\infty$.
  \item Ambillah $b_n = n$: pulihkan teorema rata-rata Cesàro pada
        \cref{exo:b1:seq:10}. Lalu tunjukkan bahwa konvers
        Cesàro--Stolz gagal: karena untuk $a_n = (-1)^n$, $b_n = n$,
        hasil bagi $a_n/b_n$ konvergen sedangkan hasil bagi
        selisihnya tidak. Jadi Stolz jalan satu arah.
\end{enumerate}

\medskip
\noindent\textbf{Bagian III --- Dividen yang pertama.}
\begin{enumerate}[resume]
  \item Buktikan $(1+h)^{3/2} - 1 = \dfrac{3h + 3h^2 +
        h^3}{(1+h)^{3/2} + 1}$ lewat pensekawanan, lalu simpulkan
        $n\bigl((1 + \tfrac1n)^{3/2} - 1\bigr) \to \tfrac32$,
        dan tutuplah dengan Cesàro--Stolz:
        \[
        T_n = \sum_{k=1}^{n} \sqrt k \;\sim\; \tfrac23\,
        n^{3/2} ,
        \]
        yang menuntaskan gantungan pada pertanyaan 2.
  \item Dari (G1) semata, turunkan apitan logaritmanya
        \[
        \frac{t}{1 + t} \;\leq\; \ln(1 + t) \;\leq\; t
        \qquad (t > -1)
        \]
        \emph{(terapkan (G1) pada $u = \ln(1+t)$ dan pada $u =
        -t/(1+t)$)}.
  \item Tunjukkan bahwa $b_n = \ln n$ naik tegas dengan
        $\ln n \to +\infty$, lalu buktikan dengan Cesàro--Stolz dan
        pertanyaan 9 bahwa
        \[
        H_n = \sum_{k=1}^{n} \frac 1k \;\sim\; \ln n .
        \]
        (Adapun struktur yang lebih halus $H_n = \ln n + \gamma + o(1)$ merupakan
        soal akhir pekan \cref{ch:b1:series}.)
  \item (Dari nisbah ke akar) Misalkan $u_n > 0$ dengan
        $\frac{u_{n+1}}{u_n} \to L > 0$. Dengan memakai pertanyaan 9, tunjukkan
        $\ln\frac{u_{n+1}}{u_n} \to \ln L$; lalu terapkan Cesàro untuk
        menyimpulkan $\frac{\ln u_n}{n} \to \ln L$, kemudian, dengan (G1),
        bahwa $u_n^{1/n} \to L$. Penerapannya: hitunglah
        $\lim\,\binom{2n}{n}^{1/n}$.
\end{enumerate}

\medskip
\noindent\textbf{Bagian IV --- Jatuhnya sinus yang lambat.} Misalkan
$u_0 \in \R$ dan $u_{n+1} = \sin u_n$.
\begin{enumerate}[resume]
  \item Dari (G2), tunjukkan $0 < \sin x < x$ untuk $0 < x \leq 1$.
        Simpulkan: bahwa $u_1 \in \intcc{-1}{1}$; bahwa jika $u_1 = 0$ maka
        barisannya nol mulai peringkat $1$; dan bahwa jika $u_1 > 0$ (dengan
        kasus $u_1 < 0$ yang simetris, karena $\sin$ ganjil), maka
        $(u_n)_{n \geq 1}$ turun tegas, positif, dan
        konvergen ke $0$ \emph{(kenalilah limitnya lewat $\ell =
        \sin \ell$, dengan memakai $\abs{\sin a - \sin b} \leq \abs{a -
        b}$, yang sendirinya merupakan akibat (G2) dan rumus
        hasil kali menjadi jumlah)}.
  \item Anggaplah mulai sekarang $u_1 \in \intoc{0}{1}$. Tunjukkan lewat
        pengapitan, dengan memakai (G2):
        \[
        \frac{\sin u_n}{u_n} \to 1
        \qquad\text{dan}\qquad
        \frac{u_n - \sin u_n}{u_n^{3}} \to \frac16 .
        \]
  \item Buktikan pemfaktoran
        \[
        w_n := \frac{1}{u_{n+1}^{2}} - \frac{1}{u_n^{2}}
        = \frac{u_n - \sin u_n}{u_n^{3}} \cdot
        \frac{u_n + \sin u_n}{u_n} \cdot
        \Bigl(\frac{u_n}{\sin u_n}\Bigr)^{2},
        \]
        lalu simpulkan $w_n \to \frac13$.
  \item Tutuplah dengan \cref{exo:b1:seq:10} (yaitu versi selisihnya)
        bahwa $\frac{1}{n\,u_n^{2}} \to \frac13$, lalu, lewat
        argumen pensekawanan bagi akar kuadratnya, judulnya:
        \[
        \sqrt n\;u_n \longrightarrow \sqrt 3 ,
        \qquad\text{yakni}\qquad
        u_n \sim \sqrt{\frac 3n} .
        \]
  \item Kuantifikasikan kelambanannya: tunjukkan bahwa akhirnya $\sqrt{2/n}
        \leq u_n \leq 2/\sqrt n$, sehingga mencapai $u_n \leq
        10^{-2}$ menuntut lebih dari $20\,000$ iterasi (kira-kira
        $30\,000$, menurut asimtotiknya). Bandingkanlah dengan \omterm{ex:b1:seq:heron}{metode
        Heron} (\cref{ex:b1:seq:heron}) lalu jelaskan
        alasan strukturalnya: bahwa pada titik tetap $0$, kemiringan
        $\sin$ adalah $1$ (yaitu titik tetap yang \emph{netral}), sedangkan
        iterasi yang memaruhkan galat menuntut kemiringan bermodulus $< 1$.
  \item Tunjukkan bahwa untuk \emph{setiap} titik awal $u_0 \in \R$,
        entah $u_n = 0$ mulai peringkat $1$ ke depan, entah $\abs{u_n} \sim
        \sqrt{3/n}$ --- jadi jatuhnya bersifat semesta, dan hanya tandanya
        yang mengingat $u_0$.
\end{enumerate}

\medskip
\noindent\textbf{Bagian V --- Asas umumnya.} Sinus itu
satu contoh dari sebuah mesin.
\begin{enumerate}[resume]
  \item Misalkan $u_n > 0$, $u_n \to 0$, dan $\dfrac{u_n -
        u_{n+1}}{u_n^{2}} \to a > 0$. Buktikan berturut-turut:
        $\frac{u_{n+1}}{u_n} \to 1$; lalu $\frac{1}{u_{n+1}} -
        \frac{1}{u_n} \to a$; lalu $n\,u_n \to \frac1a$.
  \item (Model yang persis) Untuk $u_{n+1} = \dfrac{u_n}{1 + u_n}$, $u_0
        > 0$: tunjukkan bahwa $\frac{1}{u_n}$ bersifat aritmetika, pecahkan dengan persis,
        lalu periksa kesimpulan pertanyaan 18 terhadap rumus yang persis itu.
  \item Untuk $u_{n+1} = u_n \eu^{-u_n}$, $u_0 > 0$: tunjukkan $u_n \to
        0$, pakailah (G1) untuk mengapit $\frac{1 - \eu^{-t}}{t}$ di antara
        $\frac{1}{1+t}$ dan $1$ untuk $t > 0$, lalu simpulkan $u_n
        \sim \frac 1n$.
  \item (Kontak kubik, teleskop yang dikuadratkan) Misalkan $u_n > 0$, $u_n
        \to 0$, $\dfrac{u_n - u_{n+1}}{u_n^{3}} \to a > 0$.
        Sesuaikan pemfaktoran pertanyaan 14 untuk menunjukkan
        $\frac{1}{u_{n+1}^2} - \frac{1}{u_n^2} \to 2a$, lalu
        simpulkan $n\,u_n^{2} \to \frac{1}{2a}$. Periksa bahwa $a =
        \frac16$ memulihkan Bagian IV.
\end{enumerate}

\medskip
\noindent\textbf{Bagian VI --- Batas metodenya, dan moralnya.}
\begin{enumerate}[resume]
  \item Tunjukkan bahwa hipotesis $b_n \to +\infty$ tak dapat
        dilepaskan: karena untuk $a_n = 2 - 2^{-n}$ dan $b_n = 1 - 2^{-n}$,
        hasil bagi selisihnya menuju $1$ sedangkan $\frac{a_n}
        {b_n} \to 2$. Tunjuklah baris yang persis pada bukti pertanyaan
        5 yang patah.
  \item (Stolz dua kali) Buktikan $\sum_{k=1}^{n} H_k \sim n \ln n$
        \emph{(dengan satu penerapan Cesàro--Stolz, lalu pertanyaan
        10; batasilah $(n+1)\ln(n+1) - n\ln n$ memakai pertanyaan 9)}.
  \item (Rata-rata geometri) Jika $u_n > 0$ dan $u_n \to \ell > 0$,
        tunjukkan $(u_1 u_2 \cdots u_n)^{1/n} \to \ell$; sedangkan jika $u_n \to
        +\infty$, tunjukkan $(u_1 \cdots u_n)^{1/n} \to +\infty$.
        Simpulkan $(n!)^{1/n} \to +\infty$.
  \item Sintesis, satu kalimat untuk masing-masing: (i) di manakah persisnya
        kelengkapan masuk ke soal ini; (ii) dalam arti apakah
        Cesàro--Stolz merupakan aturan l'Hospital yang diskret (adapun kembaran
        diferensialnya bersandar pada teorema nilai rata-rata pada
        \cref{ch:b1:derivative}); (iii) nyatakanlah heuristik yang
        menghubungkan orde kontak $f$ pada titik tetap yang
        netral dengan eksponen peluruhan $u_{n+1} = f(u_n)$;
        (iv) telusurilah konstanta $3$ pada $\sqrt{3/n}$ mundur lewat
        pipa $\frac16 \to \frac13 \to 3$.

\end{enumerate}
\end{problem}
